\documentclass[12pt,a4paper]{article}
\usepackage[utf8]{inputenc}
\usepackage[T1]{fontenc}
\usepackage[english]{babel}
\usepackage{geometry}
\usepackage{graphicx}
\usepackage{booktabs}
\usepackage{listings}
\usepackage{xcolor}
\usepackage{hyperref}
\usepackage{fancyhdr}
\usepackage{amsmath}
\usepackage{array}
\usepackage{tabularx}
\usepackage{float}
\usepackage{parskip}
\usepackage{mdframed}
\usepackage{enumitem}
\usepackage{titlesec}

\geometry{margin=2.5cm}

% Colors
\definecolor{codeblue}{RGB}{0,100,180}
\definecolor{codegray}{RGB}{245,245,245}
\definecolor{codegreen}{RGB}{0,128,0}
\definecolor{codepurple}{RGB}{128,0,128}
\definecolor{codered}{RGB}{180,0,0}
\definecolor{primaryblue}{RGB}{26,86,219}
\definecolor{lightblue}{RGB}{235,245,255}

% Listings style
\lstdefinestyle{sqlstyle}{
  language=SQL,
  backgroundcolor=\color{codegray},
  basicstyle=\ttfamily\small,
  keywordstyle=\color{codeblue}\bfseries,
  commentstyle=\color{codegreen}\itshape,
  stringstyle=\color{codepurple},
  breaklines=true,
  frame=single,
  rulecolor=\color{gray!40},
  numbers=none,
  captionpos=b,
  showstringspaces=false,
  tabsize=2
}

\lstdefinestyle{bashstyle}{
  language=bash,
  backgroundcolor=\color{black!90},
  basicstyle=\ttfamily\small\color{white},
  keywordstyle=\color{yellow},
  commentstyle=\color{green!60},
  stringstyle=\color{orange},
  breaklines=true,
  frame=single,
  numbers=none,
  captionpos=b,
  showstringspaces=false
}

% Header/Footer
\pagestyle{fancy}
\fancyhf{}
\fancyhead[L]{\small Distributed Systems -- Workshop 5}
\fancyhead[R]{\small Semester II 2026}
\fancyfoot[C]{\thepage}
\renewcommand{\headrulewidth}{0.4pt}

% Section formatting
\titleformat{\section}{\large\bfseries\color{primaryblue}}{}{0em}{\thesection\quad}[\titlerule]
\titleformat{\subsection}{\normalsize\bfseries}{}{0em}{\thesubsection\quad}

\hypersetup{
  colorlinks=true,
  linkcolor=primaryblue,
  urlcolor=codeblue,
  citecolor=codeblue
}

%------------------------------------------------------------
\begin{document}

%------------------------------------------------------------
\begin{titlepage}
  \centering
  \vspace*{1cm}
  {\Huge\bfseries\color{primaryblue} Distributed Databases\\[0.4em]}
  {\large\color{gray} Replication, Sharding, and Concurrency Control\\[2em]}
  \rule{\linewidth}{0.5pt}\\[1.5em]
  {\Large Workshop 5 -- Distributed Systems\\[0.5em]}
  {\large Semester II, 2026\\[0.5em]}
  {\large Prof.\ Francisco Hidrobo\\[2em]}
  \rule{\linewidth}{0.5pt}\\[2em]
  \begin{tabular}{ll}
    \textbf{Student:} & Dario Pomasqui \\[1.5em]
    \textbf{Date:} & October 1, 2026 \\
    \textbf{Repository:} & \href{https://github.com/adrianasofiaespinoza/Sistemas-Distribuidos/tree/main/Workshop5}{adrianasofiaespinoza/Sistemas-Distribuidos} \\
  \end{tabular}
\end{titlepage}


%------------------------------------------------------------
\section{Docker Topology -- Activity 1}

\subsection{Infrastructure Overview}

The workshop provisions three independent MySQL~8.0 containers using Docker
Compose, forming two distinct distributed-database patterns:

\begin{itemize}[leftmargin=1.5cm]
  \item \textbf{mysql-primary} (port 3306) -- The primary write node.
        Binary logging and GTID mode are enabled so that the replica can track
        every committed transaction. This node also serves as \emph{Shard~1}
        for the horizontal-sharding exercise.
  \item \textbf{mysql-replica} (port 3307) -- Read-only asynchronous replica.
        It consumes the binary log produced by the primary and applies each
        event in strict commit order.
  \item \textbf{mysql-shard2} (port 3308) -- Completely independent node that
        holds \emph{Shard~2}. It is not part of the replication chain.
\end{itemize}

\begin{figure}[H]
\centering
\begin{tabular}{|c|c|c|}
\hline
\textbf{Container} & \textbf{Host Port} & \textbf{Role} \\
\hline
mysql-primary  & 3306 & Primary / Shard 1 \\
mysql-replica  & 3307 & Replica (read-only) \\
mysql-shard2   & 3308 & Shard 2 \\
\hline
\end{tabular}
\caption{Docker container topology.}
\end{figure}

\subsection{docker-compose.yml (key excerpts)}

\begin{lstlisting}[style=sqlstyle, caption={Primary node configuration.}]
mysql-primary:
  image: mysql:8.0
  container_name: mysql-primary
  environment:
    MYSQL_ROOT_PASSWORD: root
    MYSQL_DATABASE: lab
  command:
    - --server-id=1
    - --log-bin=mysql-bin
    - --binlog-format=ROW
    - --gtid-mode=ON
    - --enforce-gtid-consistency=ON
    - --log-replica-updates=ON
  ports: ["3306:3306"]
\end{lstlisting}

\subsection{Verification}

After running \texttt{docker compose up -d}, all three nodes responded to
\texttt{mysqladmin ping}:

\begin{lstlisting}[style=bashstyle, caption={All nodes alive.}]
$ docker exec mysql-primary mysqladmin ping -uroot -proot
mysqld is alive
$ docker exec mysql-replica mysqladmin ping -uroot -proot
mysqld is alive
$ docker exec mysql-shard2  mysqladmin ping -uroot -proot
mysqld is alive
\end{lstlisting}

%------------------------------------------------------------
\section{Primary--Replica Replication -- Activity 2}

\subsection{Configuration Steps}

A dedicated replication account was created on the primary:

\begin{lstlisting}[style=sqlstyle, caption={Creating the replication user.}]
CREATE USER IF NOT EXISTS 'repl'@'%' IDENTIFIED BY 'repl';
GRANT REPLICATION SLAVE ON *.* TO 'repl'@'%';
FLUSH PRIVILEGES;
\end{lstlisting}

The replica was then pointed at the primary using GTID auto-positioning:

\begin{lstlisting}[style=sqlstyle, caption={Replica configuration.}]
STOP REPLICA;
CHANGE REPLICATION SOURCE TO
  SOURCE_HOST='mysql-primary',
  SOURCE_USER='repl',
  SOURCE_PASSWORD='repl',
  SOURCE_AUTO_POSITION=1,
  GET_SOURCE_PUBLIC_KEY=1;
START REPLICA;
\end{lstlisting}

\subsection{Replication Status Evidence}

\begin{lstlisting}[style=bashstyle, caption={SHOW REPLICA STATUS output (relevant fields).}]
Replica_IO_State: Waiting for source to send event
Source_Host:      mysql-primary
Replica_IO_Running:   Yes
Replica_SQL_Running:  Yes
Seconds_Behind_Source: 0
Retrieved_Gtid_Set: a97f157f-bdb3-11f1-b391-12c9eefbbc21:1-9
Executed_Gtid_Set:  a97f157f-bdb3-11f1-b391-12c9eefbbc21:1-9
\end{lstlisting}

Both \texttt{Replica\_IO\_Running} and \texttt{Replica\_SQL\_Running} report
\textbf{Yes}, confirming the replication channel is healthy.

\subsection{Data Replication Test}

A row was inserted on the primary and read back from the replica:

\begin{lstlisting}[style=sqlstyle, caption={Insert on primary.}]
CREATE TABLE IF NOT EXISTS replication_test
  (id INT PRIMARY KEY, note VARCHAR(80));
INSERT INTO replication_test VALUES (1, 'created on primary');
\end{lstlisting}

\begin{lstlisting}[style=bashstyle, caption={Query on replica -- row successfully replicated.}]
-- On mysql-replica:
SELECT * FROM replication_test;
+----+--------------------+
| id | note               |
+----+--------------------+
|  1 | created on primary |
+----+--------------------+
\end{lstlisting}

\subsection{Discussion Question}

\textbf{Why is the replica read-only?} \\[0.3em]
The replica is configured with \texttt{--read-only=ON} to guarantee that all
writes enter the system through a single authoritative node.  If both nodes
accepted independent writes, the binary log streams would diverge: each node
would generate its own GTID events, and those events would never be
automatically merged.  The result would be two separate histories with no
built-in reconciliation mechanism, leading to split-brain and permanent data
inconsistency.

%------------------------------------------------------------
\section{Horizontal Sharding -- Activity 3}

\subsection{Sharding Rule}

The user table is partitioned by a simple modular rule on the primary key:

\[
  \text{shard}(id) =
  \begin{cases}
    \text{Shard 1 (mysql-primary)} & \text{if } id \bmod 2 = 0 \\
    \text{Shard 2 (mysql-shard2)}  & \text{if } id \bmod 2 = 1
  \end{cases}
\]

\subsection{Schema Creation}

\begin{lstlisting}[style=sqlstyle, caption={Schema on both shards (identical structure).}]
-- Shard 1
CREATE DATABASE IF NOT EXISTS users_shard1;
CREATE TABLE IF NOT EXISTS users_shard1.users (
  id    INT PRIMARY KEY,
  name  VARCHAR(50),
  email VARCHAR(80),
  city  VARCHAR(50)
);
-- Shard 2 (same DDL, different database)
CREATE DATABASE IF NOT EXISTS users_shard2;
CREATE TABLE IF NOT EXISTS users_shard2.users ( ... );
\end{lstlisting}

\subsection{Data on Each Fragment}

\begin{lstlisting}[style=bashstyle, caption={Shard 1 -- even IDs (mysql-primary).}]
SELECT * FROM users_shard1.users ORDER BY id;
+----+-------------+----------------------+----------+
| id | name        | email                | city     |
+----+-------------+----------------------+----------+
|  2 | Ana Gomez   | ana.gomez@mail.com   | Bogota   |
|  4 | Carlos Ruiz | carlos.ruiz@mail.com | Medellin |
|  6 | Lucia Perez | lucia.perez@mail.com | Cali     |
+----+-------------+----------------------+----------+
\end{lstlisting}

\begin{lstlisting}[style=bashstyle, caption={Shard 2 -- odd IDs (mysql-shard2).}]
SELECT * FROM users_shard2.users ORDER BY id;
+----+--------------+-----------------------+--------------+
| id | name         | email                 | city         |
+----+--------------+-----------------------+--------------+
|  1 | Mario Lopez  | mario.lopez@mail.com  | Barranquilla |
|  3 | Sofia Torres | sofia.torres@mail.com | Cartagena    |
|  5 | Diego Castro | diego.castro@mail.com | Bucaramanga  |
+----+--------------+-----------------------+--------------+
\end{lstlisting}

\subsection{Reconstructed Logical Result}

Since no single SQL engine spans both servers, the coordinator (the
application layer) merges both result sets:

\begin{table}[H]
\centering
\begin{tabularx}{\textwidth}{c l l l c}
\toprule
\textbf{id} & \textbf{name} & \textbf{email} & \textbf{city} & \textbf{Shard} \\
\midrule
1 & Mario Lopez  & mario.lopez@mail.com  & Barranquilla & 2 \\
2 & Ana Gomez    & ana.gomez@mail.com    & Bogota       & 1 \\
3 & Sofia Torres & sofia.torres@mail.com & Cartagena    & 2 \\
4 & Carlos Ruiz  & carlos.ruiz@mail.com  & Medellin     & 1 \\
5 & Diego Castro & diego.castro@mail.com & Bucaramanga  & 2 \\
6 & Lucia Perez  & lucia.perez@mail.com  & Cali         & 1 \\
\bottomrule
\end{tabularx}
\caption{Full user set reconstructed by the coordinator.}
\end{table}

\subsection{Vertical Sharding Challenge}

A vertical sharding design splits columns across nodes instead of rows:

\begin{itemize}
  \item \textbf{Node A (identity):} \texttt{id}, \texttt{name}
  \item \textbf{Node B (contact):} \texttt{id}, \texttt{email}, \texttt{city}
\end{itemize}

To reconstruct a complete user, the application queries both nodes and joins
on \texttt{id}.  This approach benefits write-heavy workloads where identity
fields are updated far less often than contact details, allowing each shard to
be placed on hardware tuned for its specific access pattern.

%------------------------------------------------------------
\section{Lost Update and Locking -- Activity 4}

\subsection{4.1 -- Reproducing the Lost Update}

The accounts table was created with initial balances of 1000 for both Alice
and Bob.  Two sessions used \texttt{READ COMMITTED} isolation and performed a
naive read-modify-write cycle:

\begin{figure}[H]
\centering
\begin{tabular}{p{0.45\linewidth}|p{0.45\linewidth}}
\textbf{Session A} & \textbf{Session B} \\
\hline
\texttt{SET SESSION TRANSACTION} \newline \texttt{ISOLATION LEVEL READ COMMITTED;} &
\texttt{SET SESSION TRANSACTION} \newline \texttt{ISOLATION LEVEL READ COMMITTED;} \\
\texttt{START TRANSACTION;} & \texttt{START TRANSACTION;} \\
\texttt{SELECT balance ...} $\to$ 1000 & \texttt{SELECT balance ...} $\to$ 1000 \\
computes $1000 + 100 = 1100$ & computes $1000 - 50 = 950$ \\
\texttt{UPDATE ... SET balance=1100} & \\
\texttt{COMMIT;} & \\
& \texttt{UPDATE ... SET balance=950} \\
& \texttt{COMMIT;} \\
\end{tabular}
\caption{Lost update scenario -- Session B overwrites Session A.}
\end{figure}

\textbf{Observed final balance: 950.} The mathematically correct value, 1050,
was not preserved because Session B read the pre-A value (1000), computed its
delta from that stale baseline, and overwrote A's committed change.  Under
\texttt{READ COMMITTED}, each statement sees the latest committed snapshot, but
the application-level read and write are not atomic.

\subsection{4.2 -- Fix with SELECT \ldots FOR UPDATE}

The balance was reset to 1000.  With pessimistic locking, Session A
immediately acquires an exclusive row lock:

\begin{lstlisting}[style=sqlstyle, caption={Serialized execution with FOR UPDATE.}]
-- Session A
START TRANSACTION;
SELECT balance FROM accounts WHERE id=1 FOR UPDATE;  -- acquires lock
UPDATE accounts SET balance = balance + 100 WHERE id=1;
COMMIT;

-- Session B (blocks at FOR UPDATE until A commits)
START TRANSACTION;
SELECT balance FROM accounts WHERE id=1 FOR UPDATE;  -- waits...
UPDATE accounts SET balance = balance - 50  WHERE id=1;
COMMIT;
\end{lstlisting}

\begin{lstlisting}[style=bashstyle, caption={Final balance after locking fix.}]
SELECT id, owner, balance FROM accounts WHERE id=1;
+----+-------+---------+
| id | owner | balance |
+----+-------+---------+
|  1 | Alice |    1050 |
+----+-------+---------+
\end{lstlisting}

The final balance is \textbf{1050}, as expected.  Session B was forced to
wait for A to release the lock before reading the already-updated value (1100),
so it correctly computed $1100 - 50 = 1050$.

%------------------------------------------------------------
\section{Read-After-Write Consistency -- Activity 5}

\subsection{Asynchronous Read}

A write was performed on the primary and the replica was queried immediately:

\begin{lstlisting}[style=bashstyle, caption={Write on primary, immediate replica read.}]
-- PRIMARY
UPDATE accounts SET balance=balance+1 WHERE id=1;
SELECT id,balance,updated_at FROM accounts WHERE id=1;
+----+---------+---------------------+
| id | balance | updated_at          |
+----+---------+---------------------+
|  1 |    1051 | 2026-10-01 16:20:07 |

-- REPLICA (immediate read)
SELECT id,balance,updated_at FROM accounts WHERE id=1;
+----+---------+---------------------+
| id | balance | updated_at          |
+----+---------+---------------------+
|  1 |    1051 | 2026-10-01 16:20:07 |
\end{lstlisting}

On this fast local machine the replica had already caught up; the read was
consistent.  In production over a wide-area link, a stale value (1050) would
be observable during the replication lag window.

\subsection{GTID-Based Monotonic Read}

The primary's committed GTID set was captured:

\begin{lstlisting}[style=bashstyle, caption={GTID set on primary.}]
SELECT @@GLOBAL.gtid_executed;
a97f157f-bdb3-11f1-b391-12c9eefbbc21:1-23
\end{lstlisting}

The replica was then asked to wait until it had applied that exact set:

\begin{lstlisting}[style=sqlstyle, caption={WAIT\_FOR\_EXECUTED\_GTID\_SET on replica.}]
SELECT WAIT_FOR_EXECUTED_GTID_SET(
  'a97f157f-bdb3-11f1-b391-12c9eefbbc21:1-23', 5);
-- Returns: 0  (set fully applied within 5 s)
\end{lstlisting}

\textbf{Result: 0} -- the replica had applied the full GTID set before the
timeout.  A return value of 0 guarantees that any subsequent read on the
replica will reflect at least the writes captured in that GTID set, satisfying
a \emph{read-your-writes} requirement: after a client issues a write, it can
obtain the primary's GTID and pass it to the next read request; the read is
routed to a replica only after \texttt{WAIT\_FOR\_EXECUTED\_GTID\_SET} returns
0.

%------------------------------------------------------------
\section{Deadlock and Prevention -- Activity 6}

\subsection{6.1 -- Reproducing the Deadlock}

Two Python threads connected to the primary and acquired locks in opposite
orders:

\begin{figure}[H]
\centering
\begin{tabular}{p{0.45\linewidth}|p{0.45\linewidth}}
\textbf{Session A} & \textbf{Session B} \\
\hline
\texttt{START TRANSACTION;} & \texttt{START TRANSACTION;} \\
\texttt{UPDATE ... WHERE id=1} $\to$ locks row 1 & \\
& \texttt{UPDATE ... WHERE id=2} $\to$ locks row 2 \\
\texttt{UPDATE ... WHERE id=2} $\to$ \textbf{waits} & \\
& \texttt{UPDATE ... WHERE id=1} $\to$ \textbf{waits} \\
\textbf{COMMITTED} & \textbf{ERROR 1213 (victim)} \\
\end{tabular}
\caption{Deadlock scenario -- circular lock dependency.}
\end{figure}

\begin{lstlisting}[style=bashstyle, caption={Observed deadlock output.}]
[A] Locked row 1, waiting 2s before locking row 2...
[B] Locked row 2, waiting 2s before locking row 1...
[B] ERROR 1213: Deadlock found when trying to get lock;
                try restarting transaction
[A] COMMITTED

Session A result: COMMITTED
Session B result: ERROR 1213: Deadlock found when trying
                  to get lock; try restarting transaction
\end{lstlisting}

MySQL's deadlock detector selected Session B as the victim (it had performed
fewer row-level operations) and rolled it back automatically, releasing all of
B's locks so Session A could proceed.

\subsection{6.2 -- Prevention with Consistent Locking Order}

Both sessions were rewritten to acquire an exclusive lock on \emph{all
relevant rows in ascending id order} before performing any write:

\begin{lstlisting}[style=sqlstyle, caption={Ordered lock acquisition in both sessions.}]
START TRANSACTION;
SELECT * FROM accounts
  WHERE id IN (1,2) ORDER BY id FOR UPDATE;
-- then perform the session's specific updates
COMMIT;
\end{lstlisting}

\begin{lstlisting}[style=bashstyle, caption={No deadlock -- Session B blocks, then commits.}]
[A] Acquired locks on rows 1,2.
[B] Attempting to lock rows 1,2 in order (will block)...
[A] COMMITTED
[B] Acquired locks. Rows: [(1,'Alice',1010,...), (2,'Bob',990,...)]
[B] COMMITTED

Session A result: COMMITTED
Session B result: COMMITTED
\end{lstlisting}

\subsection{Blocking vs.\ Deadlock}

\begin{description}
  \item[Blocking] occurs when one transaction holds a lock that another
    transaction needs.  The second transaction waits until the first commits or
    rolls back.  It resolves \emph{automatically} and is an expected part of
    concurrent operation.
  \item[Deadlock] occurs when two (or more) transactions each hold a lock the
    other needs, forming a \emph{circular dependency} that can never resolve on
    its own.  The database must intervene by aborting one transaction.
\end{description}

By always acquiring multiple row locks in the same canonical order, the
circular dependency is structurally impossible: if A holds lock 1 and wants
lock 2, B -- which also wants lock 1 first -- will simply block waiting for
lock 1, breaking the cycle.

%------------------------------------------------------------
\section{Discussion}

Replication, sharding, and concurrency control each address a distinct
dimension of the distributed-database problem:

\textbf{Replication} increases \emph{availability} and \emph{read throughput}.
By maintaining a continuously updated copy of the primary's data, the system
can serve read traffic from the replica and survive a primary failure
(after a failover procedure).  The trade-off is eventual consistency:
asynchronous replication introduces a lag window during which replicas may
serve stale reads.  GTID-based mechanisms such as
\texttt{WAIT\_FOR\_EXECUTED\_GTID\_SET} allow the application to opt in to
stronger guarantees selectively.

\textbf{Fragmentation (sharding)} improves \emph{write scalability} and
\emph{storage capacity}.  Horizontal sharding distributes rows across
independent nodes, so each shard handles a fraction of the write load and
stores a fraction of the data.  Vertical sharding splits columns, enabling
separate scaling for different access patterns.  Both strategies require the
application layer (or a middleware coordinator) to route queries and merge
partial results, adding operational complexity.

\textbf{Concurrency control} preserves \emph{data integrity} under parallel
access.  Without it, interleaved read-modify-write cycles produce lost updates
(as shown in Activity 4), and unordered multi-row locking produces deadlocks
(Activity 6).  Pessimistic locking (\texttt{SELECT \ldots FOR UPDATE})
serializes conflicting transactions at the cost of reduced concurrency;
optimistic schemes (MVCC, OCC) allow more parallelism at the cost of
occasional abort-and-retry.

Together, the three mechanisms operate at complementary layers: replication
handles \emph{node-level} redundancy and read scaling; sharding handles
\emph{data-level} partitioning; and concurrency control handles
\emph{transaction-level} correctness.  A production distributed database
typically employs all three simultaneously.

%------------------------------------------------------------
\end{document}
